Falls are a leading cause of injury, loss of independence, and mortality among
the elderly, and wearable inertial sensors offer a practical means of automatic
fall detection. Because a missed fall is far more costly than a false alarm, the
detection problem is inherently recall-oriented. This thesis investigates whether
persistent homology, applied to short wearable accelerometer windows through a
delay-embedding pipeline, can support recall-oriented fall detection that is at
once rigorous in its evaluation, principled in its feature construction, and
lightweight enough for edge deployment.

The proposed pipeline transforms each three-axis recording into an
orientation-independent signal magnitude vector, reconstructs its phase space by
delay embedding in the sense of Takens, and summarises the resulting point cloud
with persistent homology in dimensions zero and one. Ten statistics of each
persistence diagram, together with four signal-level descriptors, form a compact
twenty-four-dimensional feature vector that is passed to deliberately simple
classifiers, namely logistic regression and a support vector machine. The
hyperparameters of the topological representation are selected by Bayesian
optimisation with a tree-structured Parzen estimator. The method is evaluated on
three public benchmark datasets---MobiFall, SisFall, and FAD\_40Hz---under a
subject-aware framework comprising naive, leave-one-subject-out, and
subject-personalised protocols, complemented by a decision-threshold sweep.

Under the leave-one-subject-out protocol the representation attains a mean
$F_2$ score of up to $0.97$ on SisFall, $0.96$ on FAD\_40Hz, and $0.90$ on
MobiFall, with recall above $0.95$ throughout. The difference between the two
classifiers is consistently below $0.02$, whereas tuning the topological
hyperparameters changes $F_2$ by an order of magnitude more, demonstrating that
performance is governed by the representation rather than the classifier. The
optimal topological configuration differs across datasets, reflecting genuine
differences in fall dynamics and point-cloud geometry. Post-hoc threshold
personalisation, which requires no model retraining, yields a strongly
right-skewed distribution of gains: most subjects benefit only marginally while a
small minority gain more than $0.10$ in $F_2$. The Cohen--Steiner stability
theorem provides a formal noise-robustness certificate that distinguishes the
topological representation from competing feature families.

\vspace{3mm}\noindent
\textbf{Keywords:} fall detection, wearable sensors, topological data analysis,
persistent homology, delay embedding, recall-oriented classification.
